%%%%%
%
%% Dans ce fichier, les �ventuels accents sont cod�s
%% en ISO-8859-15 (aka latin9)
%% La ligne suivante donne trois exemples 
%% e accent aigu : �, e accent grave : �, e-dans-l'o : �
\section{Indice de stationnarit� de la covariance}

\itemstyle{Op\'erateur~:} \operfmt{Stationnarite\_Covar\_XY}
\index{Stationnarite\_Covar\_XY}

%%%%%OBJET%%%%
\itemstyle{Objet~:}
Calcul d'un indice de stationnarit� de
la covariance $\overline{X'Y'}$ en comparant la covariance moyenne calcul�e sur des sous-p�riodes, et la covariance calcul�e sur la totalit� de la p�riode.

%%%%%ARGUMENTS%%%%
\bigskip
%% Dans ce fichier, les �ventuels accents sont cod�s
%% en ISO-8859-15 (aka latin9)
%% La ligne suivante donne trois exemples 
%% e accent aigu : �, e accent grave : �, e-dans-l'o : �
%%%%%ENTREES%%%%
\begin{tabarg}
\multicolumn{3}{l}{\titretab {Entr\'ees}}\\
\hline
$X$ & Vecteur, \texttt{float} & Param�tre X \\
$Y$ & Vecteur, \texttt{float} & Param�tre Y \\
$f_e$ & Coeff., \texttt{float} & Fr�quence d'�chantillonnage X et Y (Hz)\\
$T_1$ & Coeff., \texttt{int} & Dur�e de la p�riode d'int�gration partielle  (s)\\
$T_2$ & Coeff., \texttt{int} & Dur�e de la p�riode d'int�gration totale  (s)\\
$R_{manq}$ & Coeff., \texttt{int} & Fraction minimale de valeurs non manquantes tol�r�es (compris entre 0 et 1) \\
\end{tabarg}

\medskip
%%%%%SORTIES%%%%
\begin{tabarg}
\multicolumn{3}{l}{\titretab {Sorties}}\\
\hline
$I_{stat}$ & Vecteur, \texttt{float} & Indice de stationnarit� (sans unit�)\\
\end{tabarg}


%%%%%FORMULE%%%%
\itemstyle{Formule}

%% Formule de l'operateur : Stationnarite_Covar_XY
%% Dans ce fichier, les �ventuels accents sont cod�s
%% en ISO-8859-15 (aka latin9)
%% La ligne suivante donne trois exemples 
%% e accent aigu : �, e accent grave : �, e-dans-l'o : �

$M$ est le nombre de sous-p�riodes sur lesquelles on calcule la covariance. Donc $M$ vaut $T_2 /T_1$ ( $T_2$ > $T_1$ et $T_2$ multiple de $T_1$ )\\

On d�finit : \\
$N_1$ : nombre d'�chantillons dans la p�riode $T_1$ \\

La covariance sur la $i$�me p�riode $T_1$ est calcul�e par :
$$
(\overline{X'Y'})_i = \frac{1}{N_1}\left[ \sum_j X_j.Y_j - \frac{1}{N_1} \sum_j X_j \sum_j Y_j \right] 
$$
Puis la covariance moyenne vaut
$$
\overline{X'Y'}_{moy} = \frac{1}{M} \sum_i (\overline{X'Y'})_i
$$
L'indice de stationnarit� vaut :
$$
Istat = \left\vert  \frac{ \overline{ X'Y'}  - \overline{X'Y'}_{moy} }
                   { \overline{ X'Y'} } \right\vert
$$
avec $\overline{ X'Y'}$ : covariance calcul�e sur toute la p�riode $T_2$

%%%%%AUTEUR%%%%
\itemstyle{Auteur~:}
CNRM/GMEI/TRAMM

%%%%%REFERENCE%%%%
\itemstyle{R\'ef\'erences~:}
Micro-meteorology, T. Foken, 2008 \cite{MICROMETEO-FOKEN}
